\begin{figure}[ht]
\centering
\begin{tikzpicture}[scale=1.05,>=Latex]
\foreach \i in {0,...,4} {
  \foreach \j in {0,...,3} {
    \coordinate (v\i\j) at ({\i+0.5*\j},{0.8660254*\j});
    \fill[gray] (v\i\j) circle (1.1pt);
  }
}
\foreach \i in {0,...,3} {
  \foreach \j in {0,...,3} {
    \draw[gray!35] ({\i+0.5*\j},{0.8660254*\j}) --
      ({\i+1+0.5*\j},{0.8660254*\j});
  }
}
\foreach \i in {0,...,4} {
  \foreach \j in {0,...,2} {
    \draw[gray!35] ({\i+0.5*\j},{0.8660254*\j}) --
      ({\i+0.5*(\j+1)},{0.8660254*(\j+1)});
  }
}
\foreach \i in {0,...,3} {
  \foreach \j in {0,...,2} {
    \draw[gray!35] ({\i+1+0.5*\j},{0.8660254*\j}) --
      ({\i+0.5*(\j+1)},{0.8660254*(\j+1)});
  }
}
\draw[blue!65!black,very thick] (v00)--(v03);
\draw[red!65!black,very thick] (v20)--(v30);
\draw[green!45!black,thick,dashed] (v20)--(v11)--(v02)--(v03);
\draw[orange!80!black,thick,dotted] (v30)--(v21)--(v12)--(v03);
\node[below left] at (v00) {$v$};
\node[left] at (v01) {$v+e_2$};
\node[above] at (v03) {$v+3e_2$};
\node[below] at (v20) {$x_2$};
\node[below] at (v30) {$x_3$};
\draw[->] (4.5,0.35)--(5.2,0.35) node[right] {$e_1$};
\draw[->] (4.5,0.35)--(4.85,0.956) node[above right] {$e_2$};
\end{tikzpicture}
\caption{The compressed ray in the triangular exposure argument. Here
$x_k=v+ke_1$, and the highlighted test edge is $x_2x_3$. Both drawn paths
from its endpoints to $v+3e_2$ have length three. For a ray vertex
$v+\ell e_2$, the two distances are $\max(k,\ell)$ and
$\max(k+1,\ell)$, rather than
$k+\ell$ and $k+1+\ell$. The sign of
$h(v+e_2)-h(v)$ selects which facet can be exposed with uniform slack;
its near-zero case must be excluded in that argument.}
\label{fig:triangular}
\end{figure}

